%%
%% This is file `pdf-hyp.tex',
%% generated with the docstrip utility.
%%
%% The original source files were:
%%
%% pdfpages.dtx  (with options: `example2')
%% 
%% This file is a demonstration of the pdfpages package.
%% Run it through pdfLaTeX.
%% 
%% Please send error reports and suggestions for improvements to
%%   Andreas MATTHIAS <amat@kabsi.at>.
%%
\documentclass[a4paper,12pt]{article}
\usepackage[final]{pdfpages}
%% Uncomment the following lines, if you want to produce thumbnails.
%%\usepackage{pdflscape}
%%\usepackage{thumbpdf}
\usepackage{verbatim}
\usepackage[bookmarksopen]{hyperref}

\newcounter{example}
\setcounter{example}{1}

\newenvironment{example}
  {\par\vskip\topsep%
   \noindent\textbf{Example \arabic{example}:}%
   \stepcounter{example}%
   \par\vskip\topsep%
   \minipage{.9\linewidth}%
   \verbatim}
  {\endverbatim%
   \endminipage\vskip\topsep}

\newcommand{\meta}[1]{\ensuremath\langle\texttt{#1}\ensuremath\rangle}
\newcommand{\link}[2]{\meta{#1}.\meta{#2}}

\newcommand*{\bookmark}[3][0]{%
  \pdfoutline user {<< /S /GoTo /D (#3) >>} count #1 {#2}}

\pdfinfo{%
  /Title    (A Demonstration of the Hypertext Operations
             of the pdfpages Package)
  /Author   (Andreas MATTHIAS)
  /Subject  (A pdfLaTeX package to insert pages of external PDF documents.)
  /Keywords (pdfpages, pdfLaTeX, PDF, include, insert, page)
}

\begin{document}

\title{A Demonstration of the Hypertext Operations\\
       of the \texttt{pdfpages} Package}
\author{Andreas MATTHIAS}
\maketitle

This is a demonstration of the \texttt{pdfpages} package.
It is \textit{not} the documentation of the package.
To get the documentation run: `latex pdfpages.dtx'

\tableofcontents

\section{Hyperlinks}
\subsection{Links to the inserted Pages}

Hyperlinks are created by using the option \texttt{link}
of the \verb|\includepdf| command. Each inserted page
gets a link name consisting of the filename and
the page number: \link{filename}{page number}

\begin{example}
\includepdf[pages=1-2, link]{dummy.pdf}
\end{example}

In this example the two pages have the link names
`\texttt{dummy.pdf.1}' and `\texttt{dummy.pdf.2}'.
Setting links to these
pages, can be done easily with the \verb|\hyperlink|
macro from the \texttt{hyperref.sty} package:

\bigskip
\verb|\hyperlink{dummy.pdf.1}{Page 1}|\hskip10pt \hyperlink{dummy.pdf.1}{Page 1}\par
\verb|\hyperlink{dummy.pdf.2}{Page 2}|\hskip10pt \hyperlink{dummy.pdf.2}{Page 2}
\bigskip

This way you can refer to the \hyperlink{dummy.pdf.1}{first}
and the \hyperlink{dummy.pdf.2}{second page}.

\includepdf[pages=1-2, link]{dummy.pdf}
\bookmark[4]{List of Hyperfigures}{}
\bookmark[0]{Page 1}{dummy.pdf.1}
\bookmark[0]{Page 2}{dummy.pdf.2}

Inserting the same page twice would result in two
identical link names. To prevent this use, the option
\texttt{linkname} to specify another name for the links.

\begin{example}
\includepdf[pages=1-2, nup=1x2, landscape,
            link, linkname=mylink]{dummy.pdf}
\end{example}

\noindent
Now the links are called `\texttt{mylink.1}' and
`\texttt{mylink.2}'.

\bigskip
\verb|\hyperlink{mylink.1}{Page 1}|\hskip10pt \hyperlink{mylink.1}{Page 1}\par
\verb|\hyperlink{mylink.2}{Page 2}|\hskip10pt \hyperlink{mylink.2}{Page 2}

\includepdf[pages=1-2, nup=1x2, landscape,
            link, linkname=mylink]{dummy.pdf}
\bookmark[0]{Again Page 1}{mylink.1}
\bookmark[0]{Again Page 2}{mylink.2}

\subsection{Links to the original Document}

Each page can be a hyperlink to the document from
which it was extracted. This can be done with the
option \texttt{linktodoc}.

Click on the inserted pages and see what happens.

\begin{example}
\includepdf[pages=1-4, nup=2x2, linktodoc]{dummy.pdf}
\end{example}
\includepdf[pages=1-4, nup=2x2, linktodoc]{dummy.pdf}

\section{Article Threads}

In a PDF document one or more article threads may be defined.
An article thread is a logical connected sequence of content
items.

With the option \texttt{thread} the inserted pages become
an article thread.

\begin{example}
\includepdf[pages=1-4, nup=2x2, column,
            landscape, thread]{dummy-l.pdf}
\end{example}

In Acrobat Reader the mouse pointer changes to a hand with
a little arrow in it when moved over an article thread.
By clicking on the page you can easily follow the logical
structure of the article thread.

\includepdf[pages=1-4, nup=2x2, column,
            landscape, thread]{dummy-l.pdf}

\end{document}
\endinput
%%
%% End of file `pdf-hyp.tex'.
